% Deutsche Parallelfassung zu foundations/system_layers.tex

\chapter{Ausführungsreihenfolge und Verantwortung}
\label{chap:system-layers}

Die vorangegangenen Kapitel stellen bereits Begriffe, Zustand, Operatoren und
Realisierungsvertrag bereit. Dieses Kapitel zeigt ausschließlich ihre
Ausführungsreihenfolge und bestimmt, wo Verantwortung wechselt. Es führt
keinen weiteren Architekturkatalog ein.

\section{Die Kette der Krümmungsänderung}

Ausgehend von einem identifizierten Alignment und einer autorisierten
Krümmungsänderung arbeitet die formale Maschinerie wie folgt:
\[
\begin{aligned}
(\mathrm{pose2}(s_0),\widetilde\kappa)
&\xrightarrow{\ \mathcal D\ }
\widetilde x_{\mathrm{intrinsic}}
\xrightarrow{\ \mathcal G_c\ }
\widetilde x_{\mathrm{metric}}\\
&\xrightarrow{\ \mathcal P,\,\mathcal C\ }
\widetilde x_{\mathrm{comparable}}
\xrightarrow{\ \mathcal E_c\ }
y_{\mathrm{qualified}}.
\end{aligned}
\]

Die Rekonstruktion pflanzt die Krümmungsänderung fort. Metrische Realisierung
macht den Zustand messbar. Räumliche Abbildungen und
Koordinatentransformationen bringen qualifizierte Quellen in eine kompatible
Vergleichsdomäne. Die Bewertung erzeugt die Folge. Jeder Pfeil besitzt eigene
Anwendbarkeit und Provenienz; kein Pfeil ändert die Alignment-Identität.

Die Kette verwendet die Krümmungsbeziehung
\cref{eq:curvature-relation}, die Rekonstruktionsbeziehung
\cref{eq:state-evolution} und die Überhöhungsentwicklungsbeziehung
\cref{eq:cant-evolution}. Ihre Ausführung geht dem metrischen Vergleich
voraus; sie warten nicht darauf, dass eine CAD-Repräsentation das Alignment
definiert.

\section{Wo physische Realisierung und Beobachtung eintreten}

Physische Realisierung ist kein weiterer Schritt zwischen intrinsischer und
metrischer Beschreibung. Sie setzt Zustände in Beziehung:
\[
x_{\mathrm{target}}
\xrightarrow{\ \mathcal R\ }
x_{\mathrm{real}}.
\]
Beobachtung liefert über einen Messprozess Evidenz über einen Zustand. Eine
Vermessungspolyline tritt daher nicht als „Realität“ in die Hauptkette ein.
Sie tritt als qualifizierte Evidenz ein, die in eine kompatible metrische
Domäne abgebildet und mit einem qualifizierten Ziel- oder realisierten Zustand
verglichen werden kann.

\section{Verantwortung an jeder Grenze}

Die nützlichen Schichtunterscheidungen folgen aus der Operatorkette:

\begin{description}
\item[Intrinsische Verantwortung] besitzt konstruktive Steuerungen,
  Randbedingungen und Zustandsfortpflanzung.
\item[Metrische Verantwortung] besitzt metrischen Raum, Referenzfläche,
  Einbettung und Gültigkeit messbarer räumlicher Folgen.
\item[Repräsentationsverantwortung] besitzt Abtastung, Koordinatenform,
  Genauigkeit und Austauschprovenienz.
\item[Physische Verantwortung] besitzt die Beziehung zwischen Ziel- und
  realisiertem Zustand.
\item[Beobachtungsverantwortung] besitzt Quelle, Methode, Zeit, Unsicherheit
  und Evidenzgrenzen.
\item[Bewertungsverantwortung] besitzt Vergleichsverfahren, Anwendbarkeit,
  Annahmen und qualifiziertes Ergebnis---nicht die Entscheidung.
\end{description}

\begin{principle}[Prinzip der metrischen Realisierung]
\label{princ:metric-realization-layer}
Semantische Alignment-Logik wird von metrischem Realisierungsverhalten
getrennt.
\end{principle}

\begin{principle}[Prinzip der kanonischen Trennung]
\label{princ:canonical-layer-separation}
Intrinsische Geometrie, Koordinateneinbettung, metrische Realisierung,
physische Realisierung, Beobachtung und Bewertung dürfen nicht zu einer
impliziten Operation verschmolzen werden.
\end{principle}

\section{Abhängigkeiten, kein Implementierungsstack}

\begin{definition}[Kanonische Transformationsstruktur]
\label{def:canonical-transformation-structure}
Die formale Transformationsstruktur ist eine geordnete Komposition typisierter
Operatoren, deren Zwischenergebnisse Identität, Qualifikationen und Provenienz
bewahren.
\end{definition}

Die Reihenfolge beschränkt die Bedeutung, nicht das Software-Deployment.
Operationen dürfen zwischengespeichert, neu berechnet oder durch verschiedene
Werkzeuge implementiert werden, sofern ihre Verträge erhalten bleiben.

\begin{proposition}
\label{prop:alignment-not-crs-invariant}
Eine Koordinatentransformation kann unter Projektion metrische Eigenschaften
einer Repräsentation ändern; sie ändert nicht die intrinsische
Alignment-Identität.
\end{proposition}

\begin{proposition}
\label{prop:alignment-multi-layer-system}
Eine Alignment-Folge ist nur dann vertrauenswürdig, wenn die konstruktiven,
metrischen, repräsentationalen, physischen, beobachtenden und bewertenden
Verantwortlichkeiten, von denen sie abhängt, ausdrücklich bleiben.
\end{proposition}

\begin{proposition}
\label{prop:operational-behaviour-emerges}
Betriebliche Folgen entstehen aus deklarierten Kompositionen von Zustand,
Qualifikationen, Operatoren und anwendbarem Ingenieurwissen.
\end{proposition}

Das Ergebnis der Kette ist daher bescheiden, aber nützlich: Die Theorie kann
genau angeben, was sich geändert hat, welcher Zustand folgte, welche Handlung
jedes Artefakt erzeugte, wo es messbar ist, welche Freiheiten verbleiben und
welche Evidenz verglichen werden darf. Sie kann weiterhin kein berechnetes
Ergebnis in eine Ingenieurentscheidung verwandeln.

Die Ellipsoidkapitel spezialisieren diese Kette nun, indem sie einen
Erdoberflächen-Metric-Context bereitstellen.
